As Large Language Models (LLMs) continue to advance, efficient inference remains a challenge. Speculative decoding reduces sequential generation by drafting several tokens for parallel target verification. Drafters such as EAGLE and DFlash use target hidden states, making reuse harder when the target changes. We introduce RelaySpec, which adapts a pretrained DFlash drafter by training five linear maps on its target-feature inputs while keeping the draft transformer frozen. For fine-tuned targets, we train the maps to improve token proposals using target-generated responses. For cross-model transfer, we train them to reconstruct the features expected by the drafter. The maps fold into the existing fusion projection before inference. The same construction applies to EAGLE-3 with three maps, which gain 46 to 106 percent on fine-tuned targets. We evaluate adaptation to fine-tuned targets, transfer across model sizes and families, and shared serving. A Qwen3-4B drafter transferred to Qwen3-8B reaches 5.82$\times$ autoregressive throughput on MATH and 4.68$\times$ on GSM8K, recovering 100.2\% and 99.4\% of native-drafter throughput. Fitting these maps uses 16,384 training examples, about 49 times fewer than the published recipe for training a DFlash drafter, and updates no draft weight. Transfer is weaker on code and chat. A nonlinear map with more parameters does not improve on the linear one in our benchmarks. Comparisons with draft-body adaptation and feature interventions show when adapting the drafter's inputs is useful and where its benefits are limited.
